\section{Execution Flow}

The system can be used through two entry points.
\texttt{Main} starts an interactive CLI using the normal configuration from \texttt{application.conf}.
\texttt{DemoMain} starts the same actor graph with shorter delays and sends a fixed sequence of messages that exercise
the main state-machine transitions.

The normal full-arming flow is:
\begin{enumerate}
    \item the user sends \texttt{arm full PIN} through the CLI;
    \item the root actor forwards the request to the keypad actor;
    \item the keypad actor forwards the arming request to the control unit;
    \item if the PIN is correct, the control unit enters \textit{EXIT\_DELAY};
    \item when the exit timer expires, the control unit enters \textit{ARMED};
    \item a sensor activation in an active zone starts \textit{ENTRY\_DELAY};
    \item if the entry timer expires before a valid PIN, the control unit enters \textit{ALARM};
    \item submitting the correct PIN stops the siren and returns to \textit{DISARMED}.
\end{enumerate}

The partial-arming flow is the same, except that the arming request includes a selected set of zones.
For instance, night mode can arm \texttt{PERIMETER} and \texttt{GROUND\_FLOOR} while leaving
\texttt{LIVING\_AREA} and \texttt{SLEEPING\_AREA} inactive.
In that case living-room or bedroom motion events are received and logged, but they do not start the entry delay.

\medskip

\noindent\textbf{CLI-to-actor routing in \texttt{RootActor}}
\begin{lstlisting}[style=reportcode]
private Behavior<Command> onRequestPartialArming(RequestPartialArming command) {
    keypad.tell(new KeypadActor.RequestPartialArming(command.pin(), command.zones()));
    return this;
}

private Behavior<Command> onActivateFrontDoor(ActivateFrontDoor command) {
    frontDoorSensor.tell(new SensorActor.Activate());
    return this;
}
\end{lstlisting}

The CLI exposes one simulated sensor for each configured zone: front door for \texttt{PERIMETER}, ground-floor window for
\texttt{GROUND\_FLOOR}, living-room motion for \texttt{LIVING\_AREA}, and bedroom motion for
\texttt{SLEEPING\_AREA}.

The root actor does not wait synchronously for replies.
When the CLI asks for status, the root actor sends query messages to the control unit and the siren.
Their replies come back later as normal messages through Pekko message adapters.
This keeps read-only operations inside the same asynchronous message-passing model.
